\documentclass[10pt,letterpaper]{article}
\usepackage[utf8]{inputenc}
\usepackage[spanish]{babel}
\usepackage[top=1.5cm,bottom=1.5cm,left=2cm,right=2cm]{geometry}
\usepackage{tabularx}
\usepackage{xcolor}
\usepackage{hyperref}
\hypersetup{colorlinks=true, urlcolor=teal}

\newcommand{\seccion}[1]{\vspace{8pt}{\large\bfseries\color{teal}#1}\par\vspace{2pt}\hrule\vspace{6pt}}
\setlength{\parindent}{0pt}
\pagestyle{empty}

\begin{document}

{\Huge\bfseries Santiago Herrera Vélez}\\[6pt]
Ingeniero de Software \textbullet{} Medellín, Colombia\\
\href{mailto:santiago.herrera@example.com}{santiago.herrera@example.com} \textbullet{}
+57 300 555 1234 \textbullet{} \href{https://github.com/sherrera}{github.com/sherrera}

\seccion{Perfil}
Ingeniero de software con seis años de experiencia en desarrollo de sistemas
distribuidos, APIs y aplicaciones web. Interesado en arquitectura de datos y
equipos de alto rendimiento.

\seccion{Experiencia}
\begin{tabularx}{\textwidth}{@{}l X r@{}}
\textbf{Líder técnico} & Fintech Andina S.A.S. & 2023--actual\\
& \multicolumn{2}{p{0.75\textwidth}}{Diseño de la plataforma de pagos en tiempo real;
reducción de la latencia del 40\,\%. Coordinación de un equipo de ocho personas.}\\[6pt]
\textbf{Desarrollador sénior} & Logística Express & 2020--2023\\
& \multicolumn{2}{p{0.75\textwidth}}{Migración de un monolito a microservicios en Go y
Kubernetes.}\\[6pt]
\textbf{Desarrollador} & Agencia Digital Norte & 2018--2020\\
\end{tabularx}

\seccion{Educación}
\begin{tabularx}{\textwidth}{@{}X r@{}}
\textbf{Ingeniería de Sistemas}, Universidad EAFIT & 2013--2018\\
\textbf{Especialización en Analítica}, Universidad de Antioquia & 2021\\
\end{tabularx}

\seccion{Habilidades}
\begin{tabularx}{\textwidth}{@{}l X@{}}
Lenguajes: & Go, TypeScript, Python, SQL\\
Herramientas: & PostgreSQL, Kafka, Docker, Kubernetes, AWS\\
Idiomas: & Español (nativo), inglés (C1), portugués (B1)\\
\end{tabularx}

\end{document}
